\begin{table}[t]\centering\footnotesize
\setlength{\tabcolsep}{4pt}
\caption{\textbf{The cross-section helps iff the panel is cointegrated.} Causal block-gap MAE of \cafe{} vs.\ last-value (LOCF), lower is better; \emph{ratio} $=$ \cafe{}/LOCF ($<1$ $\Rightarrow$ \cafe{} wins). \emph{Top}: a controlled common-trends panel ($N{=}12$, $X_t{=}BF_t{+}U_t$, $F_t$ an $r$-dim random walk, $U_t$ stationary AR) as the cointegration rank $N{-}r$ grows $0\!\to\!9$: the ratio falls monotonically and crosses $1$. A unit-root (non-cointegrated) idiosyncratic gives $\approx1$ at every rank (control). \emph{Bottom}: two real financial panels, \emph{both} near-unit-root in levels, under the same protocol -- a cointegrated rates/yield-curve panel (one rate-level trend, stationary spreads) where \cafe{} wins $2\times$, and the independent-floating FX panel where last-value wins. The discriminating property is cointegration, not whether the data is financial.}
\label{tab:cointegration}
\begin{tabular}{@{}lrrrr@{}}
\toprule
Panel & coint.\ rank & \cafe{} & LOCF & ratio \\
\midrule
\multicolumn{5}{@{}l}{\emph{Controlled common-trends panel ($N{=}12$, block gaps)}}\\
\quad rank 0 & 0 & 0.304 & 0.217 & 1.42 \\
\quad rank 3 & 3 & 0.296 & 0.217 & 1.39 \\
\quad rank 6 & 6 & 0.225 & 0.186 & 1.23 \\
\quad rank 9 & 9 & 0.182 & 0.192 & 0.96 \\
\quad \emph{unit-root idio.\ (control)} & \multicolumn{4}{r}{no win: ratio 1.54, 1.27} \\
\midrule
\multicolumn{5}{@{}l}{\emph{Real financial panels (same protocol, both I(1))}}\\
\quad Rates / yield curve & cointeg. & \textbf{0.097} & 0.217 & \textbf{0.45} \\
\quad Exchange rates / FX & none & 0.295 & \textbf{0.100} & 2.93 \\
\bottomrule
\end{tabular}
\end{table}
